\chapter*{CAPÍTULO 6. Conclusiones, publicaciones, recomendaciones y trabajos futuros}
\addcontentsline{toc}{chapter}{CAPÍTULO 6. Conclusiones, publicaciones, recomendaciones y trabajos futuros}
\setcounter{chapter}{6}
\setcounter{section}{0}
\setcounter{subsection}{0}
\setcounter{subsubsection}{0}
\setcounter{figure}{0}
\setcounter{table}{0}
\renewcommand{\thesection}{\thechapter.\arabic{section}}
\renewcommand{\thesubsection}{\thesection.\arabic{subsection}}
\renewcommand{\thesubsubsection}{\thesubsection.\arabic{subsubsection}}
\makeatletter
\protected@edef\@currentlabel{\thechapter}%
\makeatother
\label{chap:capitulo6}

\setlength{\parindent}{0pt}

En este capítulo se presentan las conclusiones principales de la investigación, el análisis del cumplimiento de los objetivos planteados, las publicaciones derivadas del trabajo, las recomendaciones para la comunidad académica y profesional, las líneas de trabajo futuro identificadas y una reflexión final sobre el alcance y las implicaciones del modelo SocialScrum.

% ============================================================
\section{Análisis del cumplimiento de los objetivos de investigación}
\label{sec:analisis-objetivos-investigacion}
% ============================================================

A continuación se analiza el grado de cumplimiento de cada uno de los objetivos planteados en el Capítulo~\ref{chap:capitulo1}, contrastándolos con los resultados obtenidos a lo largo de la investigación.

\subsection{Objetivo general}

El objetivo general de esta investigación fue \textit{definir un proceso ágil para gestionar la deuda social generada en proyectos de desarrollo de software ágil, a partir de aspectos fundamentales propuestos en la literatura y de la adaptación de los elementos del enfoque ágil Scrum}.

Este objetivo se cumplió satisfactoriamente mediante el diseño del modelo SocialScrum, presentado en los capítulos~\ref{chap:capitulo3} y~\ref{chap:capitulo4}. SocialScrum constituye una extensión socio-técnica de Scrum que integra roles, eventos adaptados, artefactos sociales, métricas y salvaguardas éticas para la detección, análisis y mitigación de la deuda social en equipos de desarrollo de software. La validación conceptual mediante juicio de expertos (Capítulo~\ref{chap:capitulo5}) indicó que el modelo es coherente con los principios ágiles, pertinente para las necesidades de los equipos y conceptualmente viable para su implementación.

\subsection{Objetivo específico 1: Identificación de elementos fundamentales}

\textit{Identificar los elementos fundamentales para la gestión de la deuda social en el desarrollo de software ágil mediante un análisis de la literatura.}

Este objetivo se alcanzó a través de la Revisión Sistemática de la Literatura (RSL) presentada en el Capítulo~\ref{chap:capitulo2} y detallada en el Anexo~\ref{anexo:a}. La RSL permitió:

\begin{itemize}
    \item Identificar y categorizar los \gls{community-smells} más prevalentes en equipos ágiles, agrupándolos en cuatro categorías: comunicación, conocimiento, relacional y organizacional.
    \item Documentar la correlación entre los problemas sociales y las métricas técnicas negativas (aumento de defectos, reducción de velocidad, rotación de personal).
    \item Evidenciar que los marcos ágiles existentes, incluido Scrum, carecen de mecanismos explícitos para detectar y mitigar la deuda social de manera sistemática.
    \item Establecer las bases conceptuales que fundamentan la propuesta de SocialScrum, incluyendo la teoría de la autodeterminación, el modelo de confianza de Mayer, Davis y Schoorman, y los valores de Scrum reinterpretados para el ámbito social.
\end{itemize}

\subsection{Objetivo específico 2: Construcción del proceso basado en Scrum}

\textit{Construir un proceso basado en Scrum que permita gestionar la deuda social en proyectos de desarrollo de software ágil, incorporando los elementos identificados en el objetivo específico 1.}

Este objetivo se cumplió mediante el diseño e implementación teórica del modelo SocialScrum, desarrollado en dos capítulos complementarios:

\begin{itemize}
    \item En el Capítulo~\ref{chap:capitulo3} se definieron los fundamentos conceptuales del modelo: filosofía, principios, valores adaptados, modelo conceptual de roles, eventos y artefactos, así como las salvaguardas éticas y las consideraciones de escalabilidad.
    \item En el Capítulo~\ref{chap:capitulo4} se presentó la operacionalización del modelo: proceso de adopción por fases, perfil y gobernanza del Social Guide, sistema de medición Health Score, framework de priorización socio-técnica, sistema de estimación Social Story Points (SSP), eventos adaptados con protocolos detallados, artefactos sociales con plantillas operativas y una simulación de implementación que demostró la viabilidad práctica del modelo.
\end{itemize}

Las contribuciones principales del modelo incluyen: (i) el rol del Social Guide con perfil híbrido y gobernanza clara, (ii) el Health Score como métrica compuesta de cinco dimensiones, (iii) los artefactos sociales (Social Product Backlog, Registro de \gls{community-smells}, Dashboard de Métricas Sociales), (iv) el framework de priorización mediante el Triángulo de Decisión, (v) el sistema de estimación SSP, y (vi) los protocolos de confidencialidad y auditoría ética.

\subsection{Objetivo específico 3: Evaluación mediante juicio de expertos}

\textit{Evaluar el proceso propuesto mediante la técnica de validación de juicio de expertos, y determinar oportunidades de mejora a partir de un coeficiente estadístico.}

Este objetivo se cumplió a través del proceso de validación presentado en el Capítulo~\ref{chap:capitulo5}. La evaluación se realizó con la participación de 8 expertos en desarrollo de software ágil, utilizando un instrumento de 12 preguntas cerradas (escala Likert de 5 puntos) que evaluaron cuatro criterios de calidad: claridad, completitud, idoneidad y facilidad de uso. Adicionalmente, la dimensión de agilidad se evaluó mediante los 33 indicadores del framework \gls{agilitypro} \cite{ortega2019agilityevaluation}. Los principales hallazgos fueron:

\begin{itemize}
    \item Los criterios de completitud e idoneidad obtuvieron las valoraciones más altas, con consenso total ($\sigma = 0.00$) en las preguntas P5 y P9, lo que valida la solidez técnica de los artefactos sociales y la pertinencia del rol del Social Guide.
    \item La pregunta P2, relativa a la diferenciación del Social Guide respecto a otros roles, presentó la mayor discrepancia ($\sigma = 1.05$), identificándose como una oportunidad de mejora prioritaria.
    \item La matriz de trazabilidad demostró una cobertura total de los 33 indicadores de AgilityPRO, evidenciando que SocialScrum mantiene la alineación con los principios del Manifiesto Ágil.
    \item Las preguntas abiertas proporcionaron retroalimentación cualitativa sobre fortalezas (claridad estructural, alineación ágil, enfoque colaborativo) y áreas de mejora (necesidad de ejemplos prácticos, guías de adaptación y clarificación de roles).
\end{itemize}

% ============================================================
\section{Conclusiones}
\label{sec:conclusiones}
% ============================================================

A partir del trabajo realizado y los resultados obtenidos, se presentan las siguientes conclusiones:

\begin{enumerate}
    \item \textbf{La deuda social es un problema real y documentado en equipos de desarrollo de software ágil.} La RSL evidenció que los \gls{community-smells} ---como el silencio comunicativo, el aislamiento de desarrolladores y los silos organizacionales--- afectan directamente la productividad, la calidad del software y el bienestar del equipo. Sin embargo, la producción científica en este campo es limitada, lo que confirma la necesidad de propuestas como SocialScrum.

    \item \textbf{El diseño de SocialScrum sugiere que es posible integrar la gestión de la deuda social en Scrum sin comprometer la agilidad del proceso.} El modelo propuesto muestra que se pueden adaptar los eventos, artefactos y roles existentes de Scrum para incorporar una dimensión social, con una sobrecarga estimada inferior al 1\% del tiempo total del Sprint ($\sim$45 minutos por Sprint de dos semanas). La cobertura total de los 33 indicadores de AgilityPRO respalda esta alineación a nivel conceptual, aunque su confirmación práctica requiere validación empírica.

    \item \textbf{El rol del Social Guide constituye una contribución diferenciadora.} La definición de un perfil híbrido (psicología organizacional + conocimiento técnico) con gobernanza clara, salvaguardas éticas y mecanismos de auditoría permite abordar la deuda social con rigor profesional, sin solapar las funciones del Scrum Master ni del Product Owner. El consenso total de los expertos en la pregunta P9 valida la pertinencia de este rol.

    \item \textbf{Las métricas sociales propuestas ofrecen un marco para objetivar un fenómeno tradicionalmente subjetivo.} El Health Score, basado en los cinco valores de Scrum, proporciona un indicador cuantificable del estado social del equipo, con el potencial de transformar las dinámicas interpersonales en datos gestionables que alimenten las decisiones de priorización. Su validez de constructo y fiabilidad, sin embargo, deberán ser confirmadas mediante estudios psicométricos futuros.

    \item \textbf{La validación por juicio de expertos confirma la viabilidad conceptual del modelo.} Los resultados cuantitativos mostraron valoraciones predominantemente altas (escalas 4 y 5) en la mayoría de los criterios evaluados. No obstante, las discrepancias detectadas en aspectos como la diferenciación de roles (P2) señalan áreas concretas de refinamiento que deben abordarse antes de una implementación empírica.

    \item \textbf{Las salvaguardas éticas son un componente indispensable, no opcional.} La confidencialidad, el consentimiento informado y el derecho de auditoría del equipo no son recomendaciones, sino requisitos estructurales del modelo. Sin ellas, SocialScrum podría convertirse en un instrumento de vigilancia organizacional, contradiciendo su propósito fundamental.
\end{enumerate}

% ============================================================
\section{Limitaciones del trabajo}
\label{sec:limitaciones}
% ============================================================

Esta investigación presenta las siguientes limitaciones que deben considerarse al interpretar los resultados:

\begin{enumerate}
    \item \textbf{Ausencia de validación empírica:} SocialScrum no ha sido implementado en equipos de desarrollo reales. Las conclusiones sobre viabilidad se basan en análisis teórico, simulación y validación conceptual mediante juicio de expertos.

    \item \textbf{Dependencia del contexto organizacional:} El modelo asume equipos con una cultura organizacional mínimamente abierta al cambio. En contextos altamente jerárquicos, con desconfianza institucionalizada o resistentes al cambio, la adopción podría ser inviable sin un trabajo previo de transformación cultural.

    \item \textbf{Perfil del Social Guide:} La efectividad del modelo depende de encontrar personas con el perfil híbrido requerido (psicología organizacional + conocimiento técnico), lo cual puede ser difícil en mercados laborales específicos.

    \item \textbf{Instrumentos de medición no validados psicométricamente:} Las escalas del Health Score no han sido sometidas a pruebas de validez de constructo ni de confiabilidad test-retest, lo que constituye una limitación metodológica que debe abordarse en futuras iteraciones.

    \item \textbf{Sesgo cultural:} El marco se desarrolló desde una perspectiva latinoamericana y la muestra de expertos provino de contextos similares, lo que podría requerir adaptaciones para otros contextos culturales.

    \item \textbf{Tamaño de la muestra de expertos:} Aunque los 8 expertos superan el mínimo recomendado por la literatura \cite{Adillah2022}, una muestra más amplia y diversa fortalecería la generalización de los resultados.

    \item \textbf{Alcance metodológico --- Ciencia del Diseño sin intervención empírica:} Este trabajo se enmarca en un enfoque de \gls{dsr} \cite{hevner2004designscience, peffers2007dsrm}, centrado en la creación y validación conceptual del artefacto metodológico SocialScrum. Las fases ejecutadas ---identificación del problema, diseño y construcción del artefacto, evaluación mediante juicio de expertos y comunicación de resultados--- corresponden a una validación \textit{ex ante} \cite{wieringa2014dsr}. No se realizó una intervención empírica en un entorno productivo real, por lo que las conclusiones sobre la efectividad del modelo se limitan al plano conceptual. La implementación en equipos reales, mediante ciclos completos de Investigación-Acción \cite{baskerville1999actionresearch}, constituye la línea de trabajo futuro prioritaria.

    \item \textbf{Perfil disciplinar de los expertos:} El panel de validación estuvo conformado exclusivamente por profesionales de ingeniería de sistemas y telemática. Dado que SocialScrum incorpora dimensiones psicológicas (Teoría de la Autodeterminación, modelo de confianza de Mayer et al., Health Score), la ausencia de expertos en psicología organizacional, recursos humanos o ciencias sociales constituye una limitación en la validación de dichos componentes.

    \item \textbf{Métricas de consenso basadas en dispersión simple:} Para evaluar el grado de acuerdo entre los expertos se emplearon la desviación estándar ($\sigma$) y el Coeficiente de Variación (CV), los cuales son indicadores de dispersión que, si bien resultan informativos y de fácil interpretación, no capturan la estructura ordinal propia de las escalas Likert. Métricas de concordancia inter-jueces más robustas, como el coeficiente $W$ de Kendall \cite{Kendall1948} o el kappa de Fleiss \cite{Fleiss1971}, permitirían evaluar con mayor precisión la fiabilidad del consenso. Su aplicación se recomienda como trabajo futuro en iteraciones posteriores del proceso de validación.
\end{enumerate}

% ============================================================
\section{Publicaciones}
\label{sec:publicaciones}
% ============================================================

Como resultado del proceso de investigación, se generaron las siguientes producciones académicas:

\begin{enumerate}
    \item \textbf{Revisión sistemática de la literatura (RSL):} Se elaboró una RSL sobre la gestión de la deuda social en el desarrollo de software ágil, la cual sintetiza el estado del arte y fundamenta la propuesta de SocialScrum (ver Anexo~\ref{anexo:a}).

    \item \textbf{Artículo publicado en la revista AmiTIC:} Se publicó un artículo titulado \textit{``Social Debt in Agile Software Development''} en la revista AmiTIC, el cual presenta los resultados preliminares de la investigación y la propuesta del modelo SocialScrum (ver Anexo~\ref{anexo:b}).
\end{enumerate}

% ============================================================
\section{Recomendaciones}
\label{sec:recomendaciones}
% ============================================================

A partir de los resultados obtenidos y las lecciones aprendidas durante la investigación, se formulan las siguientes recomendaciones dirigidas tanto a organizaciones como a la comunidad académica.

\subsection{Para organizaciones que consideren adoptar SocialScrum}

\begin{enumerate}
    \item \textbf{Evaluar la madurez organizacional:} Antes de implementar SocialScrum, es fundamental verificar que la organización cumple con los pre-requisitos mínimos definidos en el modelo: compromiso ejecutivo genuino, cultura organizacional mínimamente saludable, autonomía del equipo y disponibilidad de un Social Guide calificado. Si estas condiciones no se cumplen, se recomienda invertir primero en un proceso de cambio cultural.

    \item \textbf{Comenzar con un piloto acotado:} Se recomienda implementar SocialScrum en un solo equipo dispuesto durante 3 a 6 Sprints antes de considerar una adopción más amplia. El Sprint~0 preparatorio es indispensable para establecer la línea base del Health Score y capacitar al equipo.

    \item \textbf{Invertir en el perfil del Social Guide:} No asignar el rol a una persona sin la formación adecuada. El perfil híbrido (psicología organizacional + conocimiento técnico) requiere una selección rigurosa y una capacitación mínima que incluya facilitación de grupos, mediación de conflictos y dominio de Scrum.

    \item \textbf{Proteger la confidencialidad desde el inicio:} Establecer de manera formal la separación entre los datos gestionados por el Social Guide y los procesos de evaluación de desempeño de recursos humanos. El Social Guide debe reportar al Scrum Master o al liderazgo técnico, nunca a recursos humanos.

    \item \textbf{Utilizar el Health Score como guía, no como objetivo:} El Health Score debe servir para orientar decisiones de priorización y detectar tendencias, no como una métrica de rendimiento que pueda ser manipulada o que genere presión artificial sobre el equipo.
\end{enumerate}

\subsection{Para la comunidad académica e investigadora}

\begin{enumerate}
    \item \textbf{Priorizar la validación empírica:} La principal recomendación para la comunidad académica es la realización de implementaciones piloto en organizaciones reales que permitan evaluar la efectividad de SocialScrum en diferentes contextos, tamaños de equipo e industrias.

    \item \textbf{Desarrollar y validar instrumentos de medición:} Se recomienda construir y validar psicométricamente las escalas del Health Score y los instrumentos de detección de \gls{community-smells}, incluyendo pruebas de validez de constructo, confiabilidad y sensibilidad al cambio.

    \item \textbf{Investigar factores contextuales:} Estudiar cómo varían la efectividad y la adopción de SocialScrum según variables como el tamaño del equipo, la modalidad de trabajo (presencial, híbrida, remota), el sector industrial y el contexto cultural.

    \item \textbf{Comparar con alternativas:} Evaluar SocialScrum frente a otras intervenciones de salud organizacional (coaching ágil, team building tradicional, programas de bienestar corporativo) para determinar su valor diferencial.

    \item \textbf{Conformar paneles de validación multidisciplinarios:} En futuras iteraciones del modelo, integrar un panel de evaluación que incluya psicólogos organizacionales, especialistas en talento humano y profesionales de ciencias sociales, además de ingenieros de software. Esto permitiría validar las dimensiones psicológicas del Health Score (autonomía, competencia, relacionalidad) y las salvaguardas éticas desde la perspectiva de las disciplinas que las fundamentan.
\end{enumerate}

% ============================================================
\section{Trabajos futuros}
\label{sec:trabajos-futuros}
% ============================================================

A partir de esta investigación, se identifican las siguientes líneas de trabajo futuro que podrían fortalecer y expandir el alcance del modelo SocialScrum:

\begin{enumerate}
    \item \textbf{Validación empírica en equipos reales mediante Investigación-Acción completa:} Ejecutar ciclos completos de Investigación-Acción (diagnóstico, acción y reflexión) \cite{baskerville1999actionresearch} en 3 a 5 equipos de desarrollo durante 6 a 12 meses, midiendo la evolución del Health Score, métricas técnicas (velocidad, defectos, rotación de personal) y satisfacción del equipo. Este paso es el complemento natural del presente trabajo, ya que permitiría pasar de la validación conceptual realizada aquí a una intervención empírica de campo que demuestre la efectividad real de SocialScrum en contextos organizacionales.

    \item \textbf{Desarrollo de herramientas tecnológicas:} Crear software de código abierto que automatice la recolección del Health Score, genere dashboards de métricas sociales y facilite el análisis de redes sociales (SNA) a partir de datos de herramientas como Slack, GitHub o Jira.

    \item \textbf{Programa de formación para Social Guides:} Diseñar y validar un programa de capacitación y certificación para Social Guides que incluya currículo estandarizado, evaluación de competencias, componente ético y prácticas supervisadas.

    \item \textbf{Integración con marcos de escalado ágil:} Explorar la compatibilidad y adaptación de SocialScrum con marcos de escalado como SAFe, LeSS o Nexus, definiendo cómo gestionar la deuda social en contextos de múltiples equipos coordinados.

    \item \textbf{Adaptación para equipos completamente remotos:} Desarrollar protocolos específicos de observación y detección de \gls{community-smells} para equipos distribuidos geográficamente, aprovechando herramientas de comunicación asíncrona y análisis de metadatos de interacción.

    \item \textbf{Automatización de detección de \gls{community-smells}:} Investigar la aplicación de técnicas de aprendizaje automático y procesamiento de lenguaje natural para detectar automáticamente patrones de deuda social a partir de datos de comunicación del equipo (mensajes, commits, revisiones de código).

    \item \textbf{Validación psicométrica del Health Score:} Someter las escalas del Health Score a un proceso formal de validación psicométrica que incluya análisis factorial, pruebas de confiabilidad (alfa de Cronbach, test-retest) y estudios de sensibilidad al cambio.

    \item \textbf{Estudio longitudinal de impacto:} Realizar un estudio longitudinal que evalúe el impacto de SocialScrum a largo plazo (12 a 24 meses) en indicadores de bienestar del equipo, retención de talento y calidad del producto de software.
\end{enumerate}

% ============================================================
\section{Reflexión final}
\label{sec:reflexion-final}
% ============================================================

La deuda social en equipos de desarrollo de software es un problema real con consecuencias medibles. Los conflictos interpersonales, la falta de comunicación, el aislamiento y las dinámicas disfuncionales no son ``problemas de recursos humanos'' desconectados de la ingeniería; son impedimentos que afectan directamente la capacidad del equipo de entregar software de calidad.

SocialScrum nació de una observación simple pero frecuentemente ignorada: las personas no son recursos intercambiables, y la calidad de sus relaciones determina, en gran medida, la calidad de lo que producen. Cuando un equipo acumula conflictos no resueltos, silencios comunicativos o conocimiento concentrado en pocas personas, está acumulando una deuda que eventualmente cobrará intereses en forma de defectos, retrasos, rotación y agotamiento profesional.

Esta investigación propuso un camino para hacer visible y gestionable esa deuda, utilizando el mismo lenguaje y las mismas estructuras que los equipos ágiles ya conocen: Sprints, backlogs, estimaciones, retrospectivas e inspección continua. No se pretendió inventar una nueva metodología, sino extender una existente para que también proteja a las personas que la ejecutan.

Los resultados de la validación por juicio de expertos son alentadores: el modelo fue percibido como claro, completo, pertinente y alineado con los principios ágiles. Sin embargo, SocialScrum aún es una propuesta conceptual que requiere el rigor de la validación empírica para demostrar su efectividad en contextos organizacionales reales. Las limitaciones identificadas ---desde la ausencia de implementación real hasta los instrumentos no validados psicométricamente--- no invalidan la propuesta, pero sí definen con claridad el camino que queda por recorrer.

El éxito de SocialScrum, como el de cualquier marco metodológico, dependerá en última instancia de la voluntad genuina de las organizaciones de invertir en el bienestar de sus equipos, no como un gesto simbólico, sino como una estrategia de sostenibilidad. Si este trabajo contribuye a que al menos un equipo de desarrollo se detenga a reflexionar sobre la salud de sus relaciones con la misma seriedad con que revisa la calidad de su código, habrá cumplido su propósito.

SocialScrum no pretende matar al Leviatán de la deuda social ---la complejidad humana es inherente al desarrollo de software y ningún proceso puede eliminarla por completo---. Lo que ofrece son las cartas de navegación y la estructura necesaria para domesticarlo: hacer visible lo que antes era invisible, gestionar lo que antes se ignoraba y proteger a las personas que, al final del día, son quienes escriben el código, diseñan las soluciones y sostienen los proyectos. Si los equipos ágiles aprenden a convivir con ese Leviatán ---a inspeccionarlo, a contenerlo, a reducir su fuerza destructiva Sprint a Sprint---, entonces no solo sobrevivirán a la tormenta de la entrega continua, sino que podrán prosperar en ella.